\chapter{Einführung}
Die Vorlesung "Einführung in Aufbau- und Verbindungstechnik und Leiterplattentechnik Praktikum" beschäftigt sich mit den Entwerfen, Bestücken, Vermessen von Leiterplatten und der tiefgreifenden Theorie. Der vorliegende Text berichtet über die praktischen Geschehnisse. 

Die weiteren Bewertungskriterien sind:

- Nachweis der KiCad-Designs (Screenshots -> es sollte erkennbar sein, dass es Ihre eigene Arbeit ist, d.h. das Namenskürzel sollte erkennbar sein)

- Kurze Beschreibung der Fertigungsprozesses der Leiterplatten

- Kurze Beschreibung der Bestückungsprozesse der Leiterplatten

- Die Messungen Ihrer Widerstands-Leiterplatten (optisch & elektrisch)

- Eine Diskussion zur Interpretation dieser Messergebnisse

- Ein Nachweis der fertig gestellten Platine(n), z.B. in Form von Fotos der fertigen Taschenlampe/Radio

- falsche Inhalte führen zu Punktabzug (z.B. falsche Ätzlösung)

\chapter(Leiterbahnwiderstand)
Die erste Aufgabe war es ein Leiterbahnwiderstand mit der Software KiCad zu designen. Mehr zu KiCad in Abschnitt 2 \ref(sec:taschenlampe).
Dabei wurden folgende Werte vorgegeben: 
\begin{itemize}
    \item 
\end{itemize}
\section(Designe)

\section(Fertigungsprozess)
\section(Messung des Leiterbahnwiderstandes)
\section(Interpretation)
\chapter(Taschenlampe)
\label(sec:taschenlampe)
\section(KiCad Library)
\section(KiCad Schematic)
\section(KiCad Layout)
\section(Bestückungsprozess)
\chapter(Radio)

